\Chapter{34}

\Verse{1}
{
	\hJ{}
}
{
	\eJ{
		34

¹ And YHWH said to Moses, “Carve two tablets of stones like the first
ones, and I’ll write on the tablets the words that were on the first
tablets, which you shattered.\\
	}
}
{
}

\Verse{2}
{
	\hJ{}
}
{
	\eJ{
		And be ready for the morning, and you shall go up in the morning to
Mount Sinai and present yourself to me there on the top of the mountain.\\
	}
}
{
}

\Verse{3}
{
	\hJ{}
}
{
	\eJ{
		And no man shall go up with you, and also let no man be seen in all of
the mountain. Also let the flock and the oxen not feed opposite that
mountain.”\\
	}
}
{
}

\Verse{4}
{
	\hJ{}
}
{
	\eJ{
		And he carved two tablets of stones like the first ones, and Moses got
up early in the morning and went up to Mount Sinai as YHWH had commanded
him, and he took in his hand two tablets of stones.\\
	}
}
{
}

\Verse{5}
{
	\hJ{}
}
{
	\eJ{
		And YHWH came down in a cloud and stood with him there, and he invoked
the name YHWH.\\
	}
}
{
}

\Verse{6}
{
	\hJ{}
}
{
	\eJ{
		And YHWH passed in front of him and called,

YHWH, YHWH, merciful and gracious God, slow to anger and abounding in
kindness and faithfulness,\\
	}
}
{
}

\Verse{7}
{
	\hJ{}
}
{
	\eJ{
		keeping kindness for thousands, bearing crime and offense and sin;
though not making one innocent: reckoning fathers’ crime on children and
on children’s children, on third generations and on fourth generations.

34:6–7. merciful, gracious, slow to anger, kindness, faithfulness,
bearing crime and offense and sin. This is possibly the most repeated
and quoted formula in the Tanak (Num 14:18–19; Jon 4:2; Joel 2:13; Mic
7:18; Pss 86:15; 103:8; 145:8; 2 Chr 30:9; Neh 9:17,31). The Torah never
says what the essence of God is, in contrast to the pagan gods. Baal is
the storm wind, Dagon is grain, Shamash is the sun. But what is YHWH?
This formula, expressed in the moment of the closest revelation any
human has of God in the Bible, is the closest the Torah comes to
describing the nature of God. Although humans are not to know what the
essence is, they can know what are the marks of the divine personality:
mercy, grace. In eight (or nine) different ways we are told of God’s
compassion. The last line of the formula (“though not making one
innocent”) conveys that this does not mean that one can just get away
with anything; there is still justice. But the formula clearly places
the weight on divine mercy over divine justice, and it never mentions
divine anger. Those who speak of the “Old Testament God of wrath” focus
disproportionately on the episodes of anger in the Bible and somehow
lose this crucial passage and the hundreds of times that the divine
mercy functions in the Hebrew Bible.\\
	}
}
{
}

\Verse{8}
{
	\hJ{}
}
{
	\eJ{
		And Moses hurried and knelt to the ground and bowed,\\
	}
}
{
}

\Verse{9}
{
	\hJ{}
}
{
	\eJ{
		and he said, “If I’ve found favor in your eyes, my Lord, may my Lord go
among us, because it is a stiff-necked people, and forgive our crime and
our sin, and make us your legacy.”\\
	}
}
{
}

\Verse{10}
{
	\hJ{}
}
{
	\eJ{
		And He said, “Here, I’m making a covenant. Before all your people I’ll
do wonders that haven’t been created in all the earth and among all the
nations; and all the people whom you’re among will see YHWH’s deeds,
because that which I’m doing with you is awesome.\\
	}
}
{
}

\Verse{11}
{
	\hJ{}
}
{
	\eJ{
		Watch yourself regarding what I command you today. Here, I’m driving
out from before you the Amorite and the Canaanite and the Hittite and
the Perizzite and the Hivite and the Jebusite.\\
	}
}
{
}

\Verse{12}
{
	\hJ{}
}
{
	\eJ{
		Be watchful of yourself that you don’t make a covenant with the
resident of the land onto which you’re coming, that he doesn’t become a
trap among you.\\
	}
}
{
}

\Verse{13}
{
	\hJ{}
}
{
	\eJ{
		But you shall demolish their altars and shatter their pillars and cut
down their Asherahs.\\
	}
}
{
}

\Verse{14}
{
	\hJ{}
}
{
	\eJ{
		For you shall not bow to another god—because YHWH: His name is
Jealous, He is a jealous God—\\
	}
}
{
}

\Verse{15}
{
	\hJ{}
}
{
	\eJ{
		that you not make a covenant with the resident of the land, and they
will prostitute themselves after their gods and sacrifice to their gods,
and he will call to you, and you will eat from his sacrifice.\\
	}
}
{
}

\Verse{16}
{
	\hJ{}
}
{
	\eJ{
		And you will take some of his daughters for your sons, and his
daughters will prostitute themselves after their gods and cause your
sons to prostitute themselves after their gods.\\
	}
}
{
}

\Verse{17}
{
	\hJ{}
}
{
	\eJ{
		You shall not make molten gods for yourself.\\
	}
}
{
}

\Verse{18}
{
	\hJ{}
}
{
	\eJ{
		You shall observe the Festival of Unleavened Bread. Seven days you
shall eat unleavened bread, which I commanded you, at the appointed
time, the month of Abib; because in the month of Abib you went out of
Egypt.\\
	}
}
{
}

\Verse{19}
{
	\hJ{}
}
{
	\eJ{
		Every first birth of a womb is mine, and all your animals that have a
male first birth, ox or sheep.\\
	}
}
{
}

\Verse{20}
{
	\hJ{}
}
{
	\eJ{
		And you shall redeem an ass’s first birth with a sheep, and if you do
not redeem it then you shall break its neck. You shall redeem every
firstborn of your sons. And none shall appear before me empty-handed.\\
	}
}
{
}

\Verse{21}
{
	\hJ{}
}
{
	\eJ{
		Six days you shall work, and in the seventh day you shall cease. In
plowing time and in harvest, you shall cease.\\
	}
}
{
}

\Verse{22}
{
	\hJ{}
}
{
	\eJ{
		And you shall make a Festival of Weeks, of the firstfruits of the
wheat harvest, and the Festival of Gathering at the end of the year.\\
	}
}
{
}

\Verse{23}
{
	\hJ{}
}
{
	\eJ{
		Three times in the year every one of your males shall appear before
the Lord YHWH, God of Israel.\\
	}
}
{
}

\Verse{24}
{
	\hJ{}
}
{
	\eJ{
		For I shall dispossess nations before you and widen your border, and
no man will covet your land while you are going up to appear before
YHWH, your God, three times in the year.\\
	}
}
{
}

\Verse{25}
{
	\hJ{}
}
{
	\eJ{
		You shall not offer the blood of my sacrifice on leavened bread.

And the sacrifice of the Festival of Passover shall not remain until the
morning.\\
	}
}
{
}

\Verse{26}
{
	\hJ{}
}
{
	\eJ{
		You shall bring the first of the firstfruits of your land to the house
of YHWH, your God.

You shall not cook a kid in its mother’s milk.\\
	}
}
{
}

\Verse{27}
{
	\hE{}
}
{
	\eE{
		And YHWH said to Moses, “Write these words for yourself, because I’ve
made a covenant with you and with Israel based on these words.”\\
	}
}
{
}

\Verse{28}
{
	\hE{}
}
{
	\eE{
		And he was there with YHWH forty days and forty nights. He did not eat
bread, and he did not drink water. And he wrote on the tablets the words
of the covenant, the Ten Commandments.\\
	}
}
{
}

\Verse{29}
{
	\hP{}
}
{
	\eP{
		And it was when Moses was coming down from Mount Sinai, and the two
tablets of the Testimony were in Moses’ hand when he was coming down
from the mountain. And Moses had not known that the skin of his face was
transformed when He was speaking with him.\\
	}
}
{
}

\Verse{30}
{
	\hP{}
}
{
	\eP{
		And Aaron and all the children of Israel saw Moses; and, here, the
skin of his face was transformed, and they were afraid of going over to
him.\\
	}
}
{
}

\Verse{31}
{
	\hP{}
}
{
	\eP{
		And Moses called to them. And Aaron and all the chiefs in the
congregation came back to him, and he spoke to them.\\
	}
}
{
}

\Verse{32}
{
	\hP{}
}
{
	\eP{
		And after that all the children of Israel went over. And he commanded
them everything that YHWH had spoken with him in Mount Sinai.\\
	}
}
{
}

\Verse{33}
{
	\hP{}
}
{
	\eP{
		And Moses finished speaking with them, and he put a veil on his face.\\
	}
}
{
}

\Verse{34}
{
	\hP{}
}
{
	\eP{
		And when Moses would come in front of YHWH to speak with Him, he would
turn away the veil until he would go out; and he would go out and speak
to the children of Israel what he had been commanded.\\
	}
}
{
}

\Verse{35}
{
	\hP{}
}
{
	\eP{
		And the children of Israel would see Moses’ face, that the skin of
Moses’ face was transformed, and Moses would put back the veil on his
face until he would come to speak with Him.\\
	}
}
{
}

